\documentclass[table]{beamer}
%[]中可以使用draft、handout、screen、transparency、trancompress、compress等参数

%指定beamer的模式与主题
\mode<presentation>
{
  \usetheme{Madrid}
%\usetheme{Boadilla}
%\usecolortheme{default}
%\usecolortheme{orchid}
%\usecolortheme{whale}
%\usefonttheme{professionalfonts}
}

%\usetheme{Madrid}
%这里还可以选择别的主题：Bergen, Boadilla, Madrid, AnnArbor, CambridgeUS, Pittsburgh, Rochester, Warsaw, ...
%有导航栏的Antibes, JuanLesPins, Montpellier, ...
%有内容的Berkeley, PaloAlto, Goettingen, Marburg, Hannover, ...
%有最小导航栏的Berlin, Ilmenau, Dresden, Darmstadt, Frankfurt, Singapore, Szeged, ...
%有章和节表单的Copenhagen, Luebeck, Malmoe, Warsaw, ...

%\usecolortheme{default}
%设置内部颜色主题(这些主题一般改变block里的颜色);这个主题一般选择动物来命名
%这里还可以选择别的颜色主题,如默认的和有特别目的的颜色主题default,structure,sidebartab,全颜色主题albatross,beetle,crane,dove,fly,seagull,wolverine,beaver

%\usecolortheme{orchid}
%设置外部颜色主题(这些主题一般改变title里的颜色);这个主题一般选择植物来命名
%这里还可以选择别的颜色主题,如默认的和有特别目的的颜色主题lily,orchid,rose

%\usecolortheme{whale}
%设置字体主题;这个主题一般选择海洋动物来命名
%这里还可以选择别的颜色主题,如默认的和有特别目的的颜色主题whale,seahorse,dolphin

%\usefonttheme{professionalfonts}
%类似的还可以定义structurebold,structuresmallcapsserif,professionalfonts

% 控制 beamer 的风格,可以根据自己的爱好修改
%\usepackage{beamerthemesplit} %使用 split 风格
%\usepackage{beamerthemeshadow} %使用 shadow 风格
%\usepackage[width=2cm,dark,tab]{beamerthemesidebar}

%插入音标
%\usepackage{tipa}
%\AtBeginDocument{
  %\renewcommand\textipa{\fontencoding{T3}\selectfont}
%}
%\AtBeginDocument{
  %\renewcommand\textipa[2][r]{{\fontfamily{cm#1}\tipaencoding #2}}
%}
%\renewenvironment{IPA}[1][r]
 %{\fontfamily{cm#1}\tipaencoding}
 %{}

% 设定英文字体
%\usepackage{fontspec}
% Fix bugs for fontspec in TeXLive2015
\ifdefined\suppressfontnotfounderror
  \expandafter\let\csname xetex_suppressfontnotfounderror:D\endcsname
    \suppressfontnotfounderror
\else
  \expandafter\let\csname xetex_suppressfontnotfounderror:D\endcsname
    \luatexsuppressfontnotfounderror
\fi
\usepackage[no-math]{fontspec}
\setmainfont{Times New Roman}
\setsansfont{Arial}
\setmonofont{Courier New}

% 设定中文字体
\usepackage[BoldFont,SlantFont,CJKchecksingle,CJKnumber]{xeCJK}
%\setCJKmainfont[BoldFont={Adobe Heiti Std},ItalicFont={Adobe Kaiti Std}]{Adobe Song Std}
\setCJKmainfont[BoldFont={Adobe Heiti Std},ItalicFont={Adobe Kaiti Std}]{WenQuanYi Micro Hei}
\setCJKsansfont{Adobe Heiti Std}
\setCJKmonofont{Adobe Fangsong Std}
\punctstyle{hangmobanjiao}

\defaultfontfeatures{Mapping=tex-text}
\usepackage{xunicode}
\usepackage{xltxtra}

\XeTeXlinebreaklocale "zh"
\XeTeXlinebreakskip = 0pt plus 1pt minus 0.1pt

\usepackage{setspace}
\usepackage{colortbl,xcolor}
\usepackage{hyperref}
%\hypersetup{xetex,bookmarksnumbered=true,bookmarksopen=true,pdfborder=1,breaklinks,colorlinks,linkcolor=blue,filecolor=black,urlcolor=cyan,citecolor=green}
\hypersetup{xetex,bookmarksnumbered=true,bookmarksopen=true,pdfborder=1,breaklinks,colorlinks,linkcolor=cyan,filecolor=black,urlcolor=blue,citecolor=green}

% 插入图片
\usepackage{graphicx}
\graphicspath{{figures/}}
% 图文混排
%\usepackage{picins}
\usepackage{floatflt}

% 可能用到的包
\usepackage{amsmath,amssymb}
%插入多媒体
%\usepackage{media9}
%\usepackage{movie15}
\usepackage{multimedia}
\usepackage{multicol}
\usepackage{multirow}

% 定义一些自选的模板，包括背景、图标、导航条和页脚等，修改要慎重
% 设置背景渐变由10%的红变成10%的结构颜色
%\beamertemplateshadingbackground{red!10}{structure!10}
%\beamertemplatesolidbackgroundcolor{white!90!blue}
% 使所有隐藏的文本完全透明、动态，而且动态的范围很小
\beamertemplatetransparentcovereddynamic
% 使itemize环境中变成小球，这是一种视觉效果
\beamertemplateballitem
% 为所有已编号的部分设置一个章节目录，并且编号显示成小球
\beamertemplatenumberedballsectiontoc
% 将每一页的要素的要素名设成加粗字体
\beamertemplateboldpartpage

% item逐步显示时,使已经出现的item、正在显示的item、将要出现的item呈现不同颜色
\def\hilite<#1>{
 \temporal<#1>{\color{gray}}{\color{blue}}
    {\color{blue!25}}
}

\renewcommand{\today}{\number\year 年 \number\month 月 \number\day 日}

%五角星
\usepackage{MnSymbol}

%去除图表标题中的figure等
\usepackage{caption}
\captionsetup{labelformat=empty,labelsep=none}

\usepackage{tabu}
\usepackage{multirow}
%表格自动换行
\usepackage{tabularx} 

% 千分号
%\usepackage{textcomp}

%罗马数字
\makeatletter
\newcommand{\rmnum}[1]{\romannumeral #1}
\newcommand{\Rmnum}[1]{\expandafter\@slowromancap\romannumeral #1@}
\makeatother

%分栏
\usepackage{multicol}

%\usepackage{enumitem}
%\usepackage{enumerate}

%键盘
\usepackage{keystroke}

%心形
%\usepackage{fdsymbol}

%插入源代码
\usepackage{listings}
\lstset{
  language=perl,                  % 程序语言名称：TeX, Perl, R, sh, bash, Awk
  basicstyle=\normalsize\tt,      %\tt指monospace字体族，程序源代码使用此族字体表示更加美观
  numbers=left,                   % 行号位置(左侧)
  numberstyle=\small,             % 行号字体的字号
  stepnumber=1,                   % 行号的显示步长
  numbersep=5pt,                  % 行号与代码间距
  backgroundcolor=\color{white},  % 背景色;需要 \usepackage{color}
  showspaces=false,               % 不显示空格
  showstringspaces=false,         % 不显示代码字符串中的空格标记
  showtabs=false,                 % 不显示 TAB
  tabsize=4, 
  frame=shadowbox,                % 把代码用带有阴影的框圈起来
  captionpos=b,                   % 标题位置
  breaklines=true,                % 对过长的代码自动断行
  breakatwhitespace=false,        % 断行只在空格处
  extendedchars=false,            % 解决代码跨页时，章节标题，页眉等汉字不显示的问题
  %escapeinside={\%*}{*},         % 跳脱字符，添加注释，暂时离开 listings 
  %escapeinside=``,
  commentstyle=\color{red!50!green!50!blue!50}\tt,  %浅灰色的注释
  rulesepcolor=\color{red!20!green!20!blue!20},     %代码块边框为淡青色
  keywordstyle=\color{blue!70}\bfseries\tt,         %代码关键字的颜色为蓝色，粗体
  identifierstyle=\tt,
  stringstyle=\tt,                % 代码字符串的特殊格式
  keepspaces=true,
  breakindent=1em,
  %breakindent=22pt,
  %breakindent=4em,
  breakautoindent=true,
  flexiblecolumns=true,
  aboveskip=1em,                  %代码块边框
  xleftmargin=2em,
  xrightmargin=2em
}

%\setbeamercolor{alerted text}{fg=magenta}
\setbeamercolor{bgcolor}{fg=yellow,bg=cyan}
%\setbeamercolor{itemize/enumerate body}{fg=green}

\begin{document}

%\includeonlyframes{current}

\logo{\includegraphics[height=0.08\textwidth]{qr.png}}

% 在每个Section前都会加入的Frame
\AtBeginSection[]
{
  \begin{frame}<beamer>
    %\frametitle{Outline}
    \frametitle{教学提纲}
    \setcounter{tocdepth}{3}
    \begin{multicols}{2}
      \tableofcontents[currentsection,currentsubsection]
      %\tableofcontents[currentsection]
    \end{multicols}
  \end{frame}
}
% 在每个Subsection前都会加入的Frame
\AtBeginSubsection[]
{
  \begin{frame}<beamer>
%%\begin{frame}<handout:0>
%% handout:0 表示只在手稿中出现
    \frametitle{教学提纲}
    \setcounter{tocdepth}{3}
    \begin{multicols}{2}
    \tableofcontents[currentsection,currentsubsection]
    \end{multicols}
%% 显示在目录中加亮的当前章节
  \end{frame}
}

% 为当前幻灯片设置背景
%{
%\usebackgroundtemplate{
%\vbox to \paperheight{\vfil\hbox to
%\paperwidth{\hfil\includegraphics[width=2in]{tijmu_charcoal.png}\hfil}\vfil}
%}
\begin{frame}[plain]
  \begin{center}
    {\Huge 疯狂的实验\\}
    \vspace{1cm}
    {\LARGE 天津医科大学\\}
    %\vspace{0.2cm}
    {\LARGE 生物医学工程与技术学院\\}
    \vspace{1cm}
    {\large 2018-2019学年下学期（春）\\ 公共选修课}
  \end{center}
\end{frame}
%}



%\includeonlyframes{current}

\title[医学]{第五章\quad 疯狂的实验——医学}
\author[Yixf]{伊现富（Yi Xianfu）}
\institute[TIJMU]{天津医科大学（TIJMU）\\ 生物医学工程与技术学院}
\date{2017年3月}

\begin{frame}
  \titlepage
\end{frame}

\begin{frame}[plain,label=current]
  \frametitle{教学提纲}
  \setcounter{tocdepth}{3}
  \begin{multicols}{2}
    \tableofcontents
  \end{multicols}
\end{frame}


\section{自体实验}
% 1-19,1802
\begin{frame}
  \frametitle{自体实验 | 黄热病}
  \begin{block}{疑问}
    黄热病能不能在人与人之间传播？（与黄热病人的接触有没有危险？）
  \end{block}
  \pause
  \begin{block}{实验}
    \begin{itemize}
      \item 1802年，斯塔宾斯\textbullet 弗斯（18岁），博士论文
      \item 狗/猫，黄热病人的呕吐物：喂食浸泡过的面包，敷在伤口上，注射到颈部静脉
      \item 他自己，呕吐物：敷在伤口上（20几个部位的重复），滴入眼睛，吸入蒸汽，吞下（制成的药片、稀释物）
      \item 他自己，其他体液或身体排泄物（血液、唾液、汗液、尿液）
    \end{itemize}
  \end{block}
  \pause
  \begin{block}{影响}
英雄之举影响甚微。实验主要揭示了黄热病无法通过一些方式传染，而人们想知道的却是：黄热病是怎么传播的？（蚊子！）
  \end{block}
\end{frame}

% 1-197,1968
\begin{frame}
  \frametitle{自体实验 | 螨虫}
  \begin{block}{现象与疑问}
    \begin{itemize}
      \item 给猫治疗耳螨时，女房东和她女儿抱怨很痒。
      \item 耳螨能够传播给人类吗？
    \end{itemize}
  \end{block}
  \vspace{-0.5em}
  \pause
  \begin{block}{实验}
    \begin{itemize}
      \item 耳螨：取自猫耳，显微镜检查以确认是耳螨
      \item 实验：将大约1克混有耳螨的猫的耳垢放到被试（洛佩兹自己）的左耳里（先后三次实验）
    \end{itemize}
  \end{block}
  \vspace{-0.5em}
  \pause
  \begin{block}{结果}
    详细记录：声响、瘙痒、疼痛、耳垢……
  \end{block}
  \vspace{-0.5em}
  \pause
  \begin{block}{启示}
    \begin{itemize}
      \item 如果一项实验无法重复，实验的结果就算不上证据确凿。
      \item 1994年凭借其论文获得搞笑诺贝尔奖！
    \end{itemize}
  \end{block}
\end{frame}

% 1-121,1946
\section{感冒}
\begin{frame}
  \frametitle{感冒 | 实验准备}
  \begin{block}{观点与现象}
    \begin{itemize}
      \item 大众观点：寒冷本身引发感冒。
      \item 去往极地进行长期考察的研究人员并不感冒。
      \item 爱斯基摩人在最冷的冬天不会生病，却在有外来船只驶入港口时生病。
    \end{itemize}
  \end{block}
  \pause
  \begin{block}{实验准备：在实验之前排除一切由外部带入的感冒可能}
    \begin{itemize}
      \item 2人一组住进套房，接下来的10天里除室友以外，要与一切没有保护装置的人保持10米以上的距离
      \item 允许散步，但不得进入医院大楼和使用交通工具
      \item 医生和护士进行检查时身穿防护服、脸配面具
      \item 每日3餐被装入保温容器，按时送到套房门前
    \end{itemize}
  \end{block}
\end{frame}

\begin{frame}
  \frametitle{感冒 | 实验}
  \begin{block}{实验}
    \begin{itemize}
      \item 将被试分成3组，每组6人
      \item 第一组接受鼻腔注入（将一位感冒患者的鼻腔分泌物过滤并稀释后滴入鼻中）
      \item 第二组接受冷却处理（先热水浸泡，后在通风的走廊里坚持带上半个小时，其余时间必须穿着湿袜子）
      \item 第三组同时施用冷却处理和鼻腔注入两种办法
    \end{itemize}
  \end{block}
  \pause
  \begin{block}{结果}
    \begin{itemize}
      \item 既受病毒感染又遭降温处理的人中有4人生病
      \item 只受病毒感染的人中有2人生病
      \item 单纯受凉并未造成任何人感冒
    \end{itemize}
  \end{block}
\end{frame}

\begin{frame}
  \frametitle{感冒 | 重复实验}
  \begin{block}{重复}
      20世纪50、60年代，上百名被试，一系列实验
  \end{block}
  \pause
  \begin{block}{结果}
    \begin{itemize}
      \item 一致：没有显示感冒与之前的受寒有何关联
      \item 不一致：只受病毒感染的小组中的感冒人数比既受感染又遭寒冷的小组中的感冒人数翻了一番
    \end{itemize}
  \end{block}
  \pause
  \begin{block}{启示}
    \begin{itemize}
      \item 科学在面对大众观念时处境艰难！
      \item 为什么感冒的冬季发病率高于夏季？（至今尚不十分明确）
      \item 冬季感冒大量发病与虚弱的免疫系统及房间的干燥空气并无关系。
      \item 更可能的情况是：冬季房间通风不良，冬季太阳照射相对较弱（紫外线可以杀灭病原）。
    \end{itemize}
  \end{block}
\end{frame}

% 
\section{}

% 
\section{}

% 
\section{}


\section{史海撷华}
% 1-4,1600
\begin{frame}
  \frametitle{史海撷华 | 称盘上的生活}
  \begin{block}{实验}
    \begin{itemize}
      \item 圣多里奥（量化实验医学的鼻祖），《静态医学医疗术》
      \item 在圣多里奥的房子里，一切都悬于称上：床、工作台、椅子……
      \item 在称盘上度过的时间之长（30年）无人能及
      \item 记录自身体重的点滴变化：进食的食物重量，排泄的废物重量，……
    \end{itemize}
  \end{block}
  \pause
  \begin{block}{论断}
    \begin{itemize}
      \item 人们所排泄的大小便仅占所进食的食品重量的很小一部分（排汗，第一个测量这一重量的人）
    \end{itemize}
  \end{block}
\end{frame}

% 
\begin{frame}
  \frametitle{史海撷华 | }
\end{frame}

% 
\begin{frame}
  \frametitle{史海撷华 | }
\end{frame}

% 
\begin{frame}
  \frametitle{史海撷华 | }
\end{frame}

% 
\begin{frame}
  \frametitle{史海撷华 | }
\end{frame}

% 
\begin{frame}
  \frametitle{史海撷华 | }
\end{frame}

% 
\begin{frame}
  \frametitle{史海撷华 | }
\end{frame}

% 
\begin{frame}
  \frametitle{史海撷华 | }
\end{frame}

% 
\begin{frame}
  \frametitle{史海撷华 | 其他}
  \begin{block}{真正的疯狂}
    \begin{itemize}
      % 1-37,1889
      \item 返老还童——注射过滤后的浆液（捣碎的睾丸+蒸馏水），安慰剂效应、荷尔蒙疗法
      % 1-47,1899
      \item 脊椎麻醉——医生助手齐上阵，频繁详细的实验记录
      \item 注射曼巴蛇毒——眼镜蛇 vs. 曼巴蛇，错误的估计、巨大的代价
      \item 
      \item 
      \item 
      \item 
      \item 
      \item 
    \end{itemize}
  \end{block}
\end{frame}



\section*{Acknowledgements}
\begin{frame}
  \frametitle{Powered by}
  \begin{center}
    \includegraphics[width=9cm]{power.png}
  \end{center}
\end{frame}

\end{document}

